\subsection{Templates}

\subsubsection{Template Parameter Kinds}

\begin{tabularx}{\linewidth}{lX}
    \hline
    Type & \texttt{typename T} or \texttt{class T}\\
         & Accepts any type as argument.\\
    \hline
    Non-type & \texttt{int N}, \texttt{size\_t N}, \texttt{auto N}, etc.\\
             & Accepts compile-time constant values.\\
             & Types: integral, pointers, references, \texttt{nullptr\_t}, enums, float (C++20).\\
    \hline
    Template template & \texttt{template <typename> class C}\\
                      & \texttt{template <typename> typename C} (C++17+)\\
                      & Accepts template as argument.\\
    \hline
\end{tabularx}

\subsubsection{Template Parameter Syntax}

\begin{tabularx}{\linewidth}{lX}
    \hline
    Basic & \texttt{typename T}\\
          & \texttt{int N}\\
          & \texttt{template <typename> class C}\\
    \hline
    With default & \texttt{typename T = int}\\
                 & \texttt{int N = 10}\\
                 & \texttt{template <typename> class C = std::vector}\\
    \hline
    Variadic & \texttt{typename... Args}\\
             & \texttt{int... Ns}\\
             & \texttt{template <typename> class... Containers}\\
    \hline
\end{tabularx}

\begin{mdframed}[nobreak=true,backgroundcolor=magenta!10,linecolor=magenta]
\textcolor{magenta}{\textbf{NOTE:}} \texttt{typename} vs \texttt{class} for type parameters are interchangeable. \texttt{typename} is preferred (more explicit). Also, \texttt{typename} has another use: disambiguating dependent types. Eg: \texttt{typename T::value\_type x;}
\end{mdframed}

\subsubsection{Template Specialization}

\textbf{Full/Explicit Specialization}\\
Provide custom implementation for specific type(s).

Eg:
\begin{mdframed}[nobreak=true,backgroundcolor=gray!10,linecolor=Firebrick4]
\begin{verbatim}
// Primary template
template <typename T>
class MyClass { /*...*/ };

// Full specialization for bool
template <>
class MyClass<bool> { /*...*/ };
\end{verbatim}
\end{mdframed}

\textbf{Partial Specialization}\\
Specialize for a subset of types. \textcolor{red}{Only for class templates, not functions.}

Eg:
\begin{mdframed}[nobreak=true,backgroundcolor=gray!10,linecolor=Firebrick4]
\begin{verbatim}
// Primary template
template <typename T, typename U>
class Pair { /*...*/ };

// Partial specialization: both types same
template <typename T>
class Pair<T, T> { /*...*/ };

// Partial specialization: pointer types
template <typename T, typename U>
class Pair<T*, U*> { /*...*/ };
\end{verbatim}
\end{mdframed}

\begin{tabularx}{\linewidth}{lX}
    \hline
    Full specialization & \texttt{template <> class C<int> \{...\}}\\
                        & All parameters specified.\\
                        & Works for classes and functions.\\
    \hline
    Partial specialization & \texttt{template <typename T> class C<T*> \{...\}}\\
                           & Some parameters remain generic.\\
                           & \textcolor{red}{Only for class templates.}\\
    \hline
\end{tabularx}

\subsubsection{Parameter Packs}

\textbf{Syntax}\\
\begin{tabularx}{\linewidth}{lX}
    \hline
    Declaration & \texttt{typename... Args}\\
                & \texttt{int... Ns}\\
    \hline
    Expansion & \texttt{Args...} expands to \texttt{T1, T2, T3, ...}\\
              & \texttt{func(args)...} expands to \texttt{func(a1), func(a2), ...}\\
              & \texttt{sizeof...(Args)} gives pack size.\\
    \hline
\end{tabularx}

\textbf{Example: Compile-time Fibonacci sequence}

\begin{mdframed}[nobreak=true,backgroundcolor=gray!10,linecolor=Firebrick4]
\begin{verbatim}
// Fibonacci calculation
template <int N>
struct Fib {
    static constexpr int value =
        Fib<N-1>::value + Fib<N-2>::value;
};
template <>
struct Fib<0> { static constexpr int value = 0; };
template <> 
struct Fib<1> { static constexpr int value = 1; };

// Variadic template to hold sequence
template <int... Ns>
struct FibSeq {
    static constexpr int values[] = { Fib<Ns>::value... };
};

// Usage
constexpr int fibs[] = FibSeq<0, 1, 2, 3, 4, 5>::values;
// fibs = {0, 1, 1, 2, 3, 5}
\end{verbatim}
\end{mdframed}

\textcolor{blue}{Pack expansion} \texttt{Fib<Ns>::value...} expands to\\
\texttt{Fib<0>::value, Fib<1>::value, ..., Fib<5>::value}

